% Options for packages loaded elsewhere
\PassOptionsToPackage{unicode}{hyperref}
\PassOptionsToPackage{hyphens}{url}
\documentclass[
  11pt,
]{article}
\usepackage{xcolor}
\usepackage[margin=1in]{geometry}
\usepackage{amsmath,amssymb}
\setcounter{secnumdepth}{-\maxdimen} % remove section numbering
\usepackage{iftex}
\ifPDFTeX
  \usepackage[T1]{fontenc}
  \usepackage[utf8]{inputenc}
  \usepackage{textcomp} % provide euro and other symbols
\else % if luatex or xetex
  \usepackage{unicode-math} % this also loads fontspec
  \defaultfontfeatures{Scale=MatchLowercase}
  \defaultfontfeatures[\rmfamily]{Ligatures=TeX,Scale=1}
\fi
\usepackage{lmodern}
\ifPDFTeX\else
  % xetex/luatex font selection
\fi
% Use upquote if available, for straight quotes in verbatim environments
\IfFileExists{upquote.sty}{\usepackage{upquote}}{}
\IfFileExists{microtype.sty}{% use microtype if available
  \usepackage[]{microtype}
  \UseMicrotypeSet[protrusion]{basicmath} % disable protrusion for tt fonts
}{}
\makeatletter
\@ifundefined{KOMAClassName}{% if non-KOMA class
  \IfFileExists{parskip.sty}{%
    \usepackage{parskip}
  }{% else
    \setlength{\parindent}{0pt}
    \setlength{\parskip}{6pt plus 2pt minus 1pt}}
}{% if KOMA class
  \KOMAoptions{parskip=half}}
\makeatother
\usepackage{longtable,booktabs,array}
\newcounter{none} % for unnumbered tables
\usepackage{calc} % for calculating minipage widths
% Correct order of tables after \paragraph or \subparagraph
\usepackage{etoolbox}
\makeatletter
\patchcmd\longtable{\par}{\if@noskipsec\mbox{}\fi\par}{}{}
\makeatother
% Allow footnotes in longtable head/foot
\IfFileExists{footnotehyper.sty}{\usepackage{footnotehyper}}{\usepackage{footnote}}
\makesavenoteenv{longtable}
\setlength{\emergencystretch}{3em} % prevent overfull lines
\providecommand{\tightlist}{%
  \setlength{\itemsep}{0pt}\setlength{\parskip}{0pt}}
\usepackage{amssymb}
\usepackage{amsmath}
\usepackage{newunicodechar}
\newunicodechar{ℂ}{\ensuremath{\mathbb{C}}}
\newunicodechar{ℝ}{\ensuremath{\mathbb{R}}}
\newunicodechar{ℍ}{\ensuremath{\mathbb{H}}}
\newunicodechar{ℓ}{\ensuremath{\ell}}
\newunicodechar{→}{\ensuremath{\rightarrow}}
\newunicodechar{⟹}{\ensuremath{\Longrightarrow}}
\newunicodechar{⊂}{\ensuremath{\subset}}
\newunicodechar{∈}{\ensuremath{\in}}
\newunicodechar{Σ}{\ensuremath{\Sigma}}
\newunicodechar{π}{\ensuremath{\pi}}
\newunicodechar{τ}{\ensuremath{\tau}}
\newunicodechar{≈}{\ensuremath{\approx}}
\newunicodechar{√}{\ensuremath{\sqrt{\,}}}
\newunicodechar{−}{\ensuremath{-}}
\newunicodechar{²}{\ensuremath{^{2}}}
\newunicodechar{·}{\ensuremath{\cdot}}
\usepackage{bookmark}
\IfFileExists{xurl.sty}{\usepackage{xurl}}{} % add URL line breaks if available
\urlstyle{same}
\hypersetup{
  pdftitle={Commitment and Participation: The Minimal Ontology of Event Density},
  pdfauthor={Allen Proxmire},
  hidelinks,
  pdfcreator={LaTeX via pandoc}}

\title{Commitment and Participation: The Minimal Ontology of Event
Density}
\author{Allen Proxmire}
\date{July 2026}

\begin{document}
\maketitle

\textbf{Series:} Event Density (ED) Generative Papers ---
substrate-evaluation arc (primitive minimality / reduction)
\textbf{Status:} Publication draft (structural). A minimality audit of
the canonical 13 primitives; consolidates the primitive-reduction
program. Standalone; cold-reader accessible. \textbf{Repository target:}
\texttt{physics-papers/substrate-evaluation/} (ED-Generative) ---
substrate-evaluation

\begin{center}\rule{0.5\linewidth}{0.5pt}\end{center}

\subsection{\texorpdfstring{Preamble: What This Paper Does NOT Claim
\emph{(written first per QC discipline; abstract reconciled against
this)}}{Preamble: What This Paper Does NOT Claim (written first per QC discipline; abstract reconciled against this)}}\label{preamble-what-this-paper-does-not-claim-written-first-per-qc-discipline-abstract-reconciled-against-this}

\begin{enumerate}
\def\labelenumi{\arabic{enumi}.}
\tightlist
\item
  \textbf{The primitives are not derived from nothing.} This paper does
  not derive ED's substrate from a deeper layer; it audits the
  \emph{internal} structure of the canonical 13 (Paper\_087) and finds
  that some are facets of others. The postulational genre is unchanged.
\item
  \textbf{This is NOT a claim that ED reduces to one primitive.} The
  honest result is a \emph{pair} --- commitment and participation ---
  over a spatial arena. Commitment cannot swallow the things it acts on
  (participation, the channels, the amplitude) or the spatial stage. Any
  reading of this paper as ``ED is just the arrow'' is a misreading; the
  paper argues the opposite where the arrow overreaches.
\item
  \textbf{Three-dimensionality is not derived here.} P06 stays a
  primitive/selection. The paper reports a \emph{conditional} derivation
  candidate (via linking) and is explicit that its premise is open.
\item
  \textbf{The ``arrow as keystone'' synthesis is an account, not a
  theorem.} That commitment does many of physics' jobs is a unifying
  observation with per-job support elsewhere in the program; it is not a
  proof that the other primitives are theorems of commitment.
\item
  \textbf{Empirical values remain inherited.} The reduction is of the
  primitive \emph{structure} (which primitives are independent), not of
  the value layer (masses, couplings, cosmological parameters), which is
  inherited from measurement throughout ED.
\item
  \textbf{No new primitive or rule is introduced.} Every move is an
  analysis of the existing canonical 13; nothing is added to the
  ontology.
\end{enumerate}

\begin{center}\rule{0.5\linewidth}{0.5pt}\end{center}

\subsection{Abstract}\label{abstract}

Event Density postulates thirteen substrate primitives (Paper\_087).
This paper asks the minimality question the position paper itself flags
(§8.4): does the set contain a derivable subset? We run a structured
reduction audit --- for each primitive, is it a \textbf{facet} of
another primitive, \textbf{presupposed} by it (logically prior), or
\textbf{independent}? --- and reach a compact result. \textbf{(1)} The
primitives split cleanly into three roles: the \emph{things} commitment
acts on (participation, distinguishable channels, a complex amplitude),
the \emph{arena} it runs on (spatial homogeneity, three dimensions,
scale, kinds), and \emph{commitment itself} with its facets
(forward-causal chains, polarity-transport, the stability landscape's
descent-\emph{direction}, and --- reclassified here --- time-homogeneity
as a process-side property). Commitment (P11) is the sole
\textbf{process}: every dynamical and directional structure in ED is
commitment acting. \textbf{(2)} One reduction and one reclassification
tighten the count. The bandwidth and phase primitives \textbf{fuse} into
one complex amplitude --- grounded on the result (inherited from the
quantum-logic reconstruction) that the amplitude field is ℂ, so
bandwidth and phase are two real readings of one complex object rather
than two primitives. And \textbf{time-homogeneity is reclassified from
the arena to the process}: canonical P13 is stated as a predicate
\emph{of} the dynamical primitive (the process-law is time-translation
invariant), so it is a property of commitment's one law, not an
independent stage-fact. This is a reclassification, not a derivation of
P13 from commitment's collapse-content --- the position paper factors
P13 out precisely because P11 does not entail it. \textbf{(3)}
Time-homogeneity does \emph{not}, however, merge with spatial
homogeneity into one spacetime symmetry: ED's time carries the arrow and
its space does not, so no unified translation symmetry exists at the
substrate level --- Poincaré invariance is emergent. The net honest
floor is \textbf{3 substrate + 4 arena + 1 commitment} (with P12's
descent-direction, but not its functional form, on the commitment side).
Its shortest \emph{naming} is the pair \textbf{commitment and
participation} --- the sole process and the primitive relation it acts
on --- offered as a mnemonic, not as a claim that the eight-item floor
collapses to two. The capstone: ED has a spatial stage but \textbf{no
separate temporal stage} --- space is a genuine arena, while time's
aspects (direction and law-invariance) are both properties of the
process. Three-dimensionality is the one arena member with a live
(conditional) route back to commitment, via a linking argument whose
premise remains open. Throughout, the reduction is of primitive
\emph{structure}, not empirical values, and it introduces no new
primitive.

\begin{center}\rule{0.5\linewidth}{0.5pt}\end{center}

\subsection{1. Introduction}\label{introduction}

Event Density is a substrate ontology built on thirteen postulated
primitives, enumerated authoritatively in the position paper
(Paper\_087). The position paper is careful to state that the set is
postulated, not derived, and it flags its own minimality as open ---
§8.4 offers a falsifier sentence: \emph{``Demonstration that the
canonical 13 contains a derivable subset (e.g., P13 derivable from P11 +
P03, or P06 derivable from P03 + P07) would reduce the primitive count
without changing downstream content.''}

This paper takes up that invitation and reports the result of a
systematic reduction audit. The question is not whether the primitives
can be derived from \emph{outside} ED (they are the foundational layer,
by construction) but whether, \emph{inside} the set, some primitives are
facets of others. The answer is yes --- at least one genuine reduction
exists (two substrate primitives fuse) --- and the shape of the answer
is more informative than the count: the primitives organize into three
roles, one of which, commitment, is (under a pivot the paper discloses)
the only source of process in the theory. It accounts for the entire
temporal side --- time's aspects turn out to be properties of the
process, not arena facts --- while leaving the spatial side, and its own
nouns, genuinely irreducible.

The result gives ED a compact, honest description. It is not ``ED is one
primitive'' (that is false, and the paper says where it fails). Its
shortest \emph{naming} is a pair --- commitment, and the participation
it acts on, over a spatial stage --- offered as a mnemonic for a fuller
eight-item floor. Section 2 states the method and its pivot caveat.
Sections 3--6 run the audit, the one reduction, and the one
reclassification. Section 7 states the floor and the capstone. Sections
8--10 give the arrow-as-keystone synthesis, limitations, and falsifiers.

\begin{center}\rule{0.5\linewidth}{0.5pt}\end{center}

\subsection{2. Method: The Reduction
Audit}\label{method-the-reduction-audit}

The primitives (Paper\_087 §5, abbreviated): P01 event-density-layer
existence; P02 participation as primitive relation; P03 channel + locus
indexing with spatial homogeneity; P04 bandwidth (non-negative additive
scalar); P05 polarity-transport along edges; P06 spatial dimension (D =
3+1); P07 distinguishable channels; P08 substrate scale ℓ\_ED; P09
polarity (U(1) phase); P10 rule-type multiplicity; P11 commitment
irreversibility; P12 stability landscape; P13 time homogeneity.

For each primitive we ask, against the canonical statement of commitment
(P11: \emph{``at commitment events, a chain's multi-channel
participation collapses to single-channel, with the un-selected
channels' phase content randomized; irreversibly''}):

\begin{itemize}
\tightlist
\item
  \textbf{FACET} --- definitionally part of, or entailed by, another
  primitive ⟹ it reduces.
\item
  \textbf{PRESUPPOSED} --- the other primitive's definition already
  requires it ⟹ it is logically \emph{prior} and cannot reduce to that
  primitive (the primitive needs it to be stated at all).
\item
  \textbf{INDEPENDENT} --- neither entailed by nor presupposed by the
  others.
\end{itemize}

The test is deliberately conservative: ``does many jobs'' is not ``is
the only primitive,'' and a primitive that P11's own definition
references cannot be a facet of P11 --- it is prior to it. The audit is
a structural analysis of the primitive statements, not an empirical
experiment; where a step rests on a measured or separately-derived
result, that is flagged in the load-bearing table (§2.5).

\textbf{A methodological caveat, stated up front.} This audit is
\emph{P11-pivoted by construction}: it tests each primitive's
reducibility \emph{relative to} commitment, not all pairwise
dependencies among the thirteen. Such a test cannot, by itself,
establish that commitment rather than some other primitive is the
natural pivot --- that choice is an \emph{input}, motivated by the
keystone pattern (§8), not an output of the audit. What the audit
\emph{can} do is show where commitment succeeds and, more importantly,
where it fails to absorb a primitive (its own nouns, its spatial stage);
those failures are genuine findings independent of the pivot choice.
Read the ``sole process'' result (§3) accordingly: it is what the
P11-pivoted audit returns, with the pivot disclosed, not a neutral proof
that no other organizing primitive is possible.

\subsubsection{2.5 Load-Bearing Step
Audit}\label{load-bearing-step-audit}

{\def\LTcaptype{none} % do not increment counter
\begin{longtable}[]{@{}
  >{\raggedright\arraybackslash}p{(\linewidth - 4\tabcolsep) * \real{0.3333}}
  >{\raggedright\arraybackslash}p{(\linewidth - 4\tabcolsep) * \real{0.3333}}
  >{\raggedright\arraybackslash}p{(\linewidth - 4\tabcolsep) * \real{0.3333}}@{}}
\toprule\noalign{}
\begin{minipage}[b]{\linewidth}\raggedright
Step
\end{minipage} & \begin{minipage}[b]{\linewidth}\raggedright
Status
\end{minipage} & \begin{minipage}[b]{\linewidth}\raggedright
Source / justification
\end{minipage} \\
\midrule\noalign{}
\endhead
\bottomrule\noalign{}
\endlastfoot
13 primitives postulated & P & Paper\_087 \\
Audit is P11-pivoted (tests reducibility \emph{relative to} commitment)
& \textbf{method caveat} & §2 --- an input motivated by the keystone
pattern, not a neutral all-pairs proof \\
Three-role split (nouns / arena / commitment+facets) &
\textbf{structural} & §3 --- audit against P11's definition \\
Commitment (P11) is the sole \emph{process} primitive &
\textbf{structural} & §3 --- dynamical/directional content is P11
acting; P01-ordering identification flagged interpretive \\
P01/P05 are commitment-activated facets & \textbf{structural} & §3 ---
forward-causality, transport-direction \\
P12: descent-\emph{direction} is a facet; functional \emph{form} (Σ
signs) is independent & \textbf{structural} & §3 --- the arrow supplies
the motion, not the law \\
Bandwidth + phase fuse into one complex amplitude & \textbf{D (grounded,
inherited)} & §4 --- grounded on the ℂ-field result (from the
reconstruction paper, as strong as it) \\
Channels (P07) and participation (P02) stay distinct &
\textbf{structural} & §4 --- basis and relation are not the amplitude
value \\
Time-homogeneity (P13) is a property of the process, not an arena
primitive & \textbf{structural (reclassification)} & §5 --- canonical
P13 is predicated \emph{of} the dynamical primitive; NOT derived from
P11's collapse-content \\
P03 and P13 do NOT merge into one spacetime homogeneity &
\textbf{structural} & §5 --- time is arrowed, space is not (rests on P11
being the sole asymmetric axis) \\
Three-dimensionality (P06) is derived & \textbf{NOT CLAIMED} & §6 ---
selected; two of three landings read 3 off observed content, one is a
conditional candidate \\
``Order is not label-held'' (3D premise input) & \textbf{measured} & §6
--- Kendall τ≈0.96 vs 0.5, certified sim (rules out labels; does not by
itself establish linking) \\
Arrow does many of physics' jobs (keystone synthesis) &
\textbf{interpretation} & §8 --- per-job support cited, not a theorem;
does not compound \\
Empirical values (masses, couplings) reduced & \textbf{NOT CLAIMED} &
preamble 5 --- value layer inherited throughout \\
\end{longtable}
}

Every step is P, method-caveat, structural, D, measured, interpretation,
or explicitly not-claimed. Note the honest tiering of the two
count-reducing moves: the amplitude fusion is \textbf{D} but
\emph{inherited} (only as strong as the ℂ-field result it rests on); the
time-homogeneity move is a \textbf{reclassification} (P13 belongs with
the process, not the arena), \textbf{not} a derivation of P13 from
commitment. Three-dimensionality is explicitly \textbf{not} claimed
derived.

\begin{center}\rule{0.5\linewidth}{0.5pt}\end{center}

\subsection{3. The Three-Role Split, and Commitment as the Sole
Process}\label{the-three-role-split-and-commitment-as-the-sole-process}

The audit produces a clean three-way partition.

\textbf{Role 1 --- the things commitment acts on (PRESUPPOSED by P11,
hence prior):} P02 participation (P11 collapses \emph{participation}),
P04 bandwidth (the collapse is bandwidth-weighted --- the Born-rule
content), P07 channels (P11 collapses \emph{to a channel}, among
\emph{distinguishable} ones), P09 phase (P11 \emph{randomizes
un-selected phases}). These four are exactly the objects P11's own
definition references; commitment cannot reduce them, it presupposes
them. They are also, together, the participation amplitude
\texttt{P\_K\ =\ √b\_K\ e\^{}\{iπ\_K\}} on channel K --- commitment is
what commits it.

\textbf{Role 2 --- the arena (INDEPENDENT of the arrow):} P03 spatial
homogeneity, P06 three dimensions, P08 scale, P10 rule-type
multiplicity, and (before §5) P13 time-homogeneity. These are the stage
the process runs on --- spatial, scale, and kind facts the arrow neither
creates nor presupposes.

\textbf{Role 3 --- commitment and its facets (the process):} P11
commitment itself; P01 chains, whose forward-directedness is the arrow
applied to participation-events; P05 polarity-transport, whose
forward-directedness is the arrow carrying the phase; and P12 the
stability landscape, but only \emph{in part} --- its
descent-\emph{direction} (that the chain moves forward down Σ) is
commitment's, while its functional \emph{form}
(\texttt{Σ\ =\ Coh\ −\ Str\ −\ Grad}, and the descent signs) is an
independent specification that commitment does not supply. So P12 is a
partial facet: the arrow provides the motion, not the law.

\textbf{One interpretive seam, flagged.} Ruling P01 (chains) a facet
identifies its \emph{forward-causal ordering} with P11's
\emph{collapse-irreversibility}. Causal precedence (a partial order on
events) and irreversibility (no map from post-state to pre-state) are,
in general physics, separable notions; ED treats them as the one arrow,
but that identification is an interpretive merge, not a forced
entailment. It is tiered structural-with-an-open-question, not settled
--- a reader who holds causal order distinct from irreversibility would
count P01's ordering as its own primitive.

The refined statement, which the rest of the paper sharpens:

\begin{quote}
\textbf{Commitment (P11) is the only PROCESS primitive.} Everything
static --- the nouns it acts on and the arena it runs on --- is not
commitment and does not reduce to it. But everything dynamical or
directional --- forward-causal chains, polarity-transport, the descent
of the stability landscape, and (as physics) measurement and the pointer
basis, matter/antimatter, gravity's causal direction, information
preservation, the hypercharge reference --- \emph{is} commitment acting.
Commitment is the single source of process; the other primitives are its
objects and its stage.
\end{quote}

This is why commitment appears everywhere in ED: not because it is the
only primitive, but because it is the only \emph{process}, so any
dynamical structure, followed to its root, is commitment.

\begin{center}\rule{0.5\linewidth}{0.5pt}\end{center}

\subsection{4. First Reduction: Bandwidth and Phase Fuse into One
Complex
Amplitude}\label{first-reduction-bandwidth-and-phase-fuse-into-one-complex-amplitude}

The four Role-1 nouns look like four primitives, but two of them fuse.
Bandwidth \texttt{b\_K\ =\ \textbar{}P\_K\textbar{}²} and phase
\texttt{π\_K\ =\ arg\ P\_K} are the modulus and phase of one complex
number \texttt{P\_K\ ∈\ ℂ}. In isolation that is only polar coordinates,
not a reduction (the complex number still has two real degrees of
freedom --- the fusion reduces the \emph{primitive count}, not the
degrees of freedom).

\textbf{What carries the reduction is one result, inherited: the
amplitude field is ℂ.} ED's amplitude field is complex --- the U(1)
phase (P09) rules out a real field, and composite structure rules out
quaternionic (this is the ℂ-selection of the quantum-logic
reconstruction paper, and the fusion is only as strong as that result,
which is itself reconstruction-tier). \emph{Because} the fundamental
object is a single ℂ-valued amplitude, bandwidth and phase are two real
\emph{readings} of that one object, not two independently-postulated
primitives. That is the 2 → 1: one complex primitive replaces a separate
real-modulus primitive and real-phase primitive.

Complementarity --- that a which-channel (bandwidth) commitment destroys
phase coherence, and a definite relative phase forces equal bandwidths
--- \emph{corroborates} this (the two readings cannot be independently
dialed), but it does not by itself do the reducing: position and
momentum are complementary yet nobody counts space as one primitive
rather than two coordinates. So the load is on the ℂ-field result;
complementarity is consistency evidence, not the criterion.

The other two Role-1 nouns stay distinct. P07 (channels) is the
\emph{basis} the amplitude lives in --- the distinguishable options, the
pointer basis --- separable from the amplitude values (channels are
distinct even at equal bandwidth and phase). P02 (participation) is the
primitive \emph{relation} --- a chain couples to a channel, here, now
--- the addressing the amplitude is content of, not the amplitude value.
So the substrate nouns compress from four to three:
\textbf{participation (a relation), channels (a basis), and the complex
amplitude (their content).}

\begin{center}\rule{0.5\linewidth}{0.5pt}\end{center}

\subsection{5. Time-Homogeneity Belongs to the Process; Space and Time
Do Not Share One
Homogeneity}\label{time-homogeneity-belongs-to-the-process-space-and-time-do-not-share-one-homogeneity}

The arena appears to have two ``no privileged spot'' members --- spatial
homogeneity (P03: no privileged locus) and time-homogeneity (P13: no
privileged moment). Two things need to be said, and they must be kept
separate because they have very different strength.

\textbf{The load-bearing, falsifiable claim: no unified spacetime
homogeneity exists at the substrate level.} The tempting move is to
merge P03 and P13 into one spacetime-translation symmetry (the
continuum's Poincaré invariance read back as primitive). ED cannot: in
ED time carries the arrow (commitment) and space does not --- there is
an arrow primitive (P11) and no spatial-arrow primitive. So
``homogeneity of space'' and ``homogeneity of time'' are homogeneities
of two structurally different axes, not one symmetry in two guises, and
a unified spacetime-translation symmetry is \emph{emergent}, not
primitive --- ED's substrate carries \textbf{less} symmetry than its
smooth limit. This is the section's real content, it is falsifiable
(§10.4), and it rests entirely on one point of failure: \textbf{P11
being the sole asymmetric axis.} If that holds, the non-merge holds.

\textbf{The reclassification (weaker, and honest about it): P13 is a
property of the process, not an arena primitive.} Here is the correction
to an earlier, circular version of this argument. One \emph{cannot}
derive P13 from P11 by observing that ``commitment is one fixed rule,''
because calling the rule ``one fixed rule'' --- the same at every moment
--- is already asserting time-homogeneity; that route is
question-begging, and the position paper factors P13 out \emph{as a
separate primitive} precisely because P11's collapse-content does not
entail it. What can be said honestly is a reclassification, and it comes
straight from the canonical text: Paper\_087 states P13 as a predicate
\textbf{of the dynamical primitive} --- \emph{``the substrate's
dynamical primitive is invariant under time translation.''} So P13 is,
by its own definition, a property of the process-law (that the one law
is time-translation invariant), \textbf{not} a fact about a separate
temporal stage the way P03 is a fact about the spatial stage. It
therefore belongs on the process side of the ledger, not in the arena
--- which removes it from the arena count (5 → 4) as a
\emph{reclassification}, not a derivation. Once it is on the process
side, the ``merge with P03'' is not even well-posed: a property of the
process cannot merge with an arena primitive. So the trivial reading of
the non-merge falls out for free; the non-trivial, falsifiable claim is
the one above.

A scope note that survives either way: P13 is about the \emph{laws}, not
the \emph{state}. The arrow may give time a special \emph{initial}
moment (a first commitment), but a special initial state is a boundary
condition, not a law-inhomogeneity --- so the arrow's possible beginning
is a separate initial-condition matter and does not bear on P13, which
is law-invariance.

\begin{center}\rule{0.5\linewidth}{0.5pt}\end{center}

\subsection{6. The One That Does Not Reduce (Yet): Three
Dimensions}\label{the-one-that-does-not-reduce-yet-three-dimensions}

Three-dimensionality (P06) is the hardest arena member, and the honest
verdict is that it is \textbf{selected, not derived} --- it stays a
primitive. Three arguments land on 3, but they are not three
derivations, and the paper should not present them as such. \textbf{Two
of the three read 3 \emph{off} observed physics rather than deriving
it:} the internal amplitude dimension (a gauge family's channels are
mutually-orthogonal complex vectors, and the max in ℂ\^{}d is d, so the
\emph{observed} families \{1,2,3\} fix d = 3 --- a read-off of the gauge
content), and the holographic reach law (only d = 3 gives the
\emph{observed} \texttt{g\ \textasciitilde{}\ 1/b} --- a read-off of
gravity). Only the third, the linking argument, is a genuine derivation
\emph{candidate}. That the three coincide on 3 is real internal
coherence, but two of them are anchored on measured content, so the
honest status is selection, not derivation.

The linking argument is the one genuine \emph{derivation candidate}, and
it routes 3D back to commitment. A link (two loops threaded) is the only
structure that holds an order against continuous rearrangement, and
stable linking of one-dimensional loops exists in exactly three
dimensions (2D cannot form a link; 4D unravels any link --- rigorous
topology). Read against the primitives, its premises are
commitment-side: ``committed order must be stable against continuous
rearrangement'' \emph{is} commitment's irreversibility (a committed
order is permanent, hence topologically protected), and ``1D loops''
\emph{are} the forward-causal chains. So the candidate is: \{1D chains +
committed order topologically stable + stable 1D-linking unique to 3D\}
⟹ d = 3 --- which, \emph{if it holds}, forces three-dimensionality from
commitment needing to hold stable topological records of 1D chains.

The candidate rests on one open premise: that ED holds its committed
order by linking \emph{specifically}. Its status is half-settled. The
cheap alternative --- order held by node labels or timestamps --- is
\textbf{ruled out} on the certified simulator: relabel the substrate
with physics fixed and the committed order barely moves (Kendall τ ≈
0.96 versus 0.5 for a label-scramble; the committing-node set is
invariant). What this measures is precisely that the order is
\textbf{not label-held}; concluding it is \emph{geometry}-held is an
inference that holds only if labels and geometry exhaust the options,
and the residual 4\% movement is not zero --- so the honest reading is
the negative one (labels ruled out), not yet the positive one. And the
structure that would hold it is reachable (the multi-chain participation
graph can be intrinsically linked). What remains open is isolating that
the holder is \emph{linking} specifically, which requires a
three-dimensional substrate to even ask, and so is not testable on the
certified two-dimensional simulator --- an honest hard limit, shared
with the open curvature-emergence derivation. So P06 stays a primitive,
with a commitment-connected, conditional route back.

\begin{center}\rule{0.5\linewidth}{0.5pt}\end{center}

\subsection{7. The Floor and the
Capstone}\label{the-floor-and-the-capstone}

Collecting the audit, the one reduction, and the one reclassification:
the substrate nouns compress 4 → 3 (§4); time-homogeneity moves from the
arena to the process side (§5); the arena's remaining members are
spatial. The honest minimal \textbf{floor} is the eight items:

\[\boxed{\;3\ \text{substrate (participation, channels, complex amplitude)} \;+\; 4\ \text{arena (spatial homogeneity, 3D, scale, kinds)} \;+\; 1\ \text{commitment (with chains, transport, and P12's descent-direction as facets; time-homogeneity as a process-side property)}\;}\]

That box is the floor. The two-word form below it is a \textbf{mnemonic
for the floor, not a replacement of it}:

\begin{quote}
\textbf{The shortest \emph{naming} of Event Density is commitment and
participation} --- the sole process, and the primitive relation it acts
on. This names the two most fundamental \emph{roles}; it is explicitly
\textbf{not} a claim that channels, the amplitude, and the four-member
spatial arena collapse into ``participation.'' Section 4 argued the
opposite (channels and amplitude stay distinct from the participation
relation), and that stands. ``Participation'' in the two-word form is
standing in, as a mnemonic, for the whole noun-side that commitment
addresses; the honest count remains the eight-item box.
\end{quote}

With that scope fixed, two features are worth stating. First, even the
mnemonic is a \emph{pair}, not a single primitive --- faithful to ED's
relational stance that a relation is prior to an isolated thing
(participation is already a relation: a chain, to a channel, here, now;
commitment is the act on that relation). Second, and the capstone:

\begin{quote}
\textbf{ED has a spatial stage but no separate temporal stage.} The
remaining arena is purely spatial --- three dimensions, spatial
homogeneity, scale, kinds --- with no arrow. Time contributes no
\emph{arena} primitive: its two aspects are commitment's direction (P11)
and the law-invariance of that one process (P13, a process-side
property, §5). There is no ``time'' sitting in the arena beside space,
because what we call time is commitment happening.
\end{quote}

This is why the two-word naming is apt and not an oversimplification:
one does not need a third word for a temporal stage, because there is no
temporal stage --- the arena is spatial, and time is the process. It
remains a \emph{naming} of roles; the floor is the box.

\begin{center}\rule{0.5\linewidth}{0.5pt}\end{center}

\subsection{8. Discussion: Commitment as the Keystone of ED's
Physics}\label{discussion-commitment-as-the-keystone-of-eds-physics}

The reduction explains a pattern visible across ED's derivations:
commitment is not merely one primitive among many but the single
load-bearing one, because it is the only process. Followed to their
roots, ED's physical structures are commitment acting:

\begin{itemize}
\tightlist
\item
  \textbf{Time and irreversibility} are commitment directly.
\item
  \textbf{Measurement and the pointer basis} --- commitment resolves
  only in the channel basis, so ED is a preferred-basis theory because
  the arrow selects the basis (einselection is primitive here).
\item
  \textbf{The measurement-repeatability (covering-law) content} ---
  irreversibility locks a committed channel, which is the repeatability
  a quantum-logic reconstruction needs.
\item
  \textbf{Matter/antimatter (C)} --- first-arrival commitment selects
  the charge conjugation, the one chirality-adjacent thing ED fixes
  natively.
\item
  \textbf{Gravity's causal direction} --- the one-way black-hole horizon
  versus the two-way cosmological one is set by the arrow.
\item
  \textbf{Information preservation} --- the same irreversibility that
  forbids reversal forbids erasure: information is preserved as the
  permanent record of every commitment.
\item
  \textbf{The single hypercharge U(1)} --- polarity is phased against
  the one external flow, the commitment flow, so there is one gauged
  U(1) because there is one arrow to phase against.
\end{itemize}

Each of these has its own support elsewhere in the program; assembled,
they are the physical face of §3's structural claim that commitment is
the sole process. This synthesis is an \textbf{account}, not a theorem
(preamble 4): it explains \emph{why} commitment is everywhere without
claiming the other primitives are derivable from it --- they are its
objects and its stage, and the paper has argued (§3--§6) exactly where
commitment stops.

One guard against a rhetorical illusion the list can create:
\textbf{each job is only as strong as its own paper, and several of the
seven are themselves interpretation-tier} (that einselection is
primitive, that matter/antimatter is native, that gravity's direction
and the single-U(1) reference are the arrow's). The list is a
\emph{pattern}, not a compounding of evidence --- seven conditional
accounts do not sum to one hard result. Read it as motivation for the
pivot choice (§2), not as independent confirmation of it.

\begin{center}\rule{0.5\linewidth}{0.5pt}\end{center}

\subsection{9. Limitations (honest)}\label{limitations-honest}

\begin{itemize}
\tightlist
\item
  \textbf{The reduction is structural, not a derivation of the
  primitives.} The primitives remain postulated (Paper\_087); this
  audits their internal dependencies. The count-reducing moves are of
  two different kinds and neither is a clean derivation: the §4
  amplitude fusion is \emph{inherited} (grounded on the ℂ-field result
  of the reconstruction paper, and only as strong as that
  reconstruction-tier result), and the §5 time-homogeneity move is a
  \emph{reclassification} (P13 belongs with the process by its canonical
  definition), explicitly \textbf{not} a derivation of P13 from
  commitment --- an earlier ``one fixed rule ⟹ time-homogeneous laws''
  version was circular and is retracted here.
\item
  \textbf{The audit is P11-pivoted} (§2). It cannot establish that
  commitment, rather than another primitive, is the natural pivot; that
  is an input. Its pivot-independent content is where commitment
  \emph{fails} to absorb a primitive.
\item
  \textbf{Three-dimensionality is not resolved.} P06 stays a primitive;
  its route back to commitment is a candidate with an open premise (§6),
  gated on the same curvature-emergence derivation that is open across
  ED's geometric arc.
\item
  \textbf{``Commitment is the sole process'' is an audit result, not a
  proof of completeness.} It rests on reading each primitive against
  P11's canonical statement; a different canonical statement of a
  primitive could shift a facet/independent call.
\item
  \textbf{The keystone synthesis is an account.} §8 assembles per-job
  results into a unifying picture; it is not a single derivation and is
  labelled interpretation.
\item
  \textbf{Values are untouched.} Nothing here reduces the empirical
  value layer; the minimality is of primitive structure only.
\end{itemize}

\begin{center}\rule{0.5\linewidth}{0.5pt}\end{center}

\subsection{10. Falsification Criteria}\label{falsification-criteria}

\subsubsection{10.1 A Role-1 noun is not presupposed by
commitment}\label{a-role-1-noun-is-not-presupposed-by-commitment}

\textbf{Falsifier sentence:} \emph{A demonstration that participation,
channels, or the amplitude can be stated without commitment presupposing
them --- or that one of them is in fact a facet of commitment --- would
overturn the three-role split (§3).}

\subsubsection{10.2 Bandwidth and phase are independent
primitives}\label{bandwidth-and-phase-are-independent-primitives}

\textbf{Falsifier sentence:} \emph{A demonstration that bandwidth and
phase are jointly sharp (no complementarity) or that the amplitude field
is not ℂ would break the 4 → 3 fusion (§4) and restore them as two
primitives.}

\subsubsection{10.3 Time-homogeneity is an arena fact, not a process
property}\label{time-homogeneity-is-an-arena-fact-not-a-process-property}

\textbf{Falsifier sentence:} \emph{A demonstration that P13
(time-homogeneity) is a fact about a temporal stage independent of the
process-law --- the way P03 is a fact about the spatial stage --- rather
than a property predicated of the dynamical primitive, would falsify the
§5 reclassification and return P13 to the arena.}

\subsubsection{10.4 Space and time share one primitive
homogeneity}\label{space-and-time-share-one-primitive-homogeneity}

\textbf{Falsifier sentence:} \emph{A demonstration that ED's substrate
carries a single unified spacetime-translation symmetry at the primitive
level (not merely emergently) would falsify the space/time asymmetry
claim (§5) and collapse the arrow's distinction between the two axes.}

\subsubsection{10.5 Three-dimensionality is genuinely derived (or
genuinely
free)}\label{three-dimensionality-is-genuinely-derived-or-genuinely-free}

\textbf{Falsifier sentence:} \emph{Either a completed derivation of P06
(closing the linking premise) or a demonstration that no consistency
argument selects 3 would move P06 off its current ``selected, not
derived'' status (§6).}

\begin{center}\rule{0.5\linewidth}{0.5pt}\end{center}

\subsection{11. Position Statement}\label{position-statement}

\textbf{The canonical thirteen primitives of Event Density organize into
three roles --- the things commitment acts on, the arena it runs on, and
commitment with its facets --- and, under a P11-pivoted audit,
commitment (P11) is the sole process primitive.} One reduction and one
reclassification tighten the set: bandwidth and phase fuse into one
complex amplitude (grounded on the inherited ℂ-field result), and
time-homogeneity is moved from the arena to the process side (P13 is, by
its canonical definition, a property of the dynamical primitive --- a
reclassification, not a derivation of P13 from commitment).
Time-homogeneity does not merge with spatial homogeneity --- ED's time
is arrowed and its space is not, so no unified spacetime symmetry exists
at the substrate level; Poincaré invariance is emergent. The honest
floor is 3 substrate + 4 arena + 1 commitment; its shortest
\emph{naming} is the pair \textbf{commitment and participation} (the
process and the relation it acts on), a mnemonic for the floor, not a
claim that the eight items collapse to two. ED has a spatial arena but
no separate temporal one: time is the process.

\textbf{What this paper claims.} The three-role split and
commitment-as-sole-process (structural, P11-pivot disclosed); the 4 → 3
amplitude fusion (grounded, inherited from the ℂ-field result); the
time-homogeneity reclassification (P13 is a process-side property, not
an arena fact); the space/time asymmetry (no unified substrate
homogeneity, resting on P11 being the sole asymmetric axis); the
eight-item floor and the spatial-stage capstone.

\textbf{What this paper does not claim.} No derivation of the primitives
from a deeper layer; no reduction to a single primitive (the two-word
form is a mnemonic, not the floor); no \emph{derivation} of
time-homogeneity from commitment (an earlier circular version is
retracted; the move is a reclassification); no derivation of
three-dimensionality (selected, with a conditional candidate); no
reduction of empirical values; and the keystone synthesis is an account,
not a theorem that does not compound.

\textbf{Series context.} The minimality companion to the position paper
(Paper\_087): where the position paper enumerates and postulates the
thirteen, this paper audits their internal dependencies and reports the
compact pair the reduction lands on. It consolidates the reduction
program (the role audit, the amplitude fusion, the three-dimensionality
test, and the time-homogeneity reclassification) into one cold-reader
statement of what Event Density is, at bottom: the pair of \emph{roles}
--- commitment and participation --- naming an eight-item floor.

\begin{center}\rule{0.5\linewidth}{0.5pt}\end{center}

\subsection{Appendix --- Glossary and Reader
Map}\label{appendix-glossary-and-reader-map}

\subsubsection{Glossary}\label{glossary}

\begin{itemize}
\tightlist
\item
  \textbf{Commitment (P11).} The irreversible collapse of a chain's
  multi-channel participation to a single channel, with un-selected
  phases randomized. ED's sole process primitive.
\item
  \textbf{Participation (P02).} The primitive relation: a chain couples
  to a channel, at a locus, at a time. The sole thing commitment acts
  on.
\item
  \textbf{Complex amplitude (P04 + P09).} \texttt{√b\ ·\ e\^{}\{iπ\}} on
  a channel --- bandwidth and phase as two real readings of one complex
  number, fused because the amplitude field is ℂ (the inherited ℂ-field
  result), with complementarity as corroboration.
\item
  \textbf{Channels (P07).} The distinguishable options / pointer basis
  the amplitude lives in; a distinct primitive from the amplitude value.
\item
  \textbf{Arena.} The stage commitment runs on: spatial homogeneity
  (P03), three dimensions (P06), scale (P08), kinds (P10). Purely
  spatial after §5.
\item
  \textbf{Facet.} A primitive whose content is definitionally part of
  another; here chains (P01) and transport (P05) are facets of
  commitment, and P12's descent-\emph{direction} is a partial facet (its
  functional form is not). Time-homogeneity (P13) is not a facet but a
  \emph{process-side property} --- reclassified out of the arena (§5),
  not derived from commitment.
\item
  \textbf{Selected, not derived.} A primitive picked out by consistency
  arguments (multiple independent ones landing on the same value)
  without being reduced to other primitives --- the status of
  three-dimensionality.
\end{itemize}

\subsubsection{Reader map}\label{reader-map}

\emph{Intuition.} Strip Event Density to its smallest honest description
and two things remain: an act (commitment) and a relation it acts on
(participation). Everything dynamical is the act; the amplitude and
channels are what the relation carries; space is the stage. Time is not
a third thing on the stage --- time is the act happening.

\emph{Where the arrow stops.} The value of this audit is as much in
where commitment \emph{fails} to reduce things as where it succeeds. It
cannot absorb its own nouns (participation, channels, amplitude) or its
spatial stage. It does absorb the entire temporal side of the arena.
That asymmetry --- space is a stage, time is the process --- is the
load-bearing, falsifiable content.

\textbf{Where to look next.} - \emph{The primitives enumerated and
postulated:} the position paper (Paper\_087). - \emph{The
complementarity and ℂ-field results the amplitude fusion rests on:} the
quantum-logic reconstruction. - \emph{The open bridge that would settle
three-dimensionality:} the curvature-emergence arc.

\textbf{Open work} (declared): close the three-dimensionality linking
premise (gated on curvature emergence); test whether participation (P02)
and channels (P07) could further merge into one relational primitive;
and formalize ``commitment is the sole process'' beyond the
primitive-statement audit.

\begin{center}\rule{0.5\linewidth}{0.5pt}\end{center}

\textbf{End of Paper (Commitment and Participation).}

\emph{Substrate-evaluation arc. A P11-pivoted minimality audit of the
canonical thirteen primitives: they split into the things commitment
acts on, the arena it runs on, and commitment with its facets;
commitment is the sole process. Bandwidth and phase fuse into one
complex amplitude (grounded on the inherited ℂ-field result);
time-homogeneity is reclassified from the arena to the process side (a
property of the one law, not derived from commitment); space and time
share no substrate homogeneity. Honest floor: 3 substrate + 4 arena + 1
commitment. Shortest naming (a mnemonic, not the floor): commitment and
participation, over a purely spatial stage --- ED has a spatial arena
but no separate temporal one, because time is the process.
Three-dimensionality stays selected, not derived, with a conditional
route back to commitment.}

\end{document}
